% IX. Live HTTP results (measured on loopback with an emulated network;
% RTT, bandwidth and object sizes are ASSUMED; constraint 8). Only valid
% passes are cited, and SSD-store passes with more than 1 % foreign I/O are
% left out (lead ruling; genlib.live_citable).
% Sources: results/live/{live_metrics,live_parity}.csv, live_config.json,
% via numbers.tex and paper/decisions.py. Generated: Table T7 (T7b and
% Fig. F7 in the appendix).

\section{Live HTTP Results (Measured)}\label{sec:live}

% t7_live.tex -- GENERATED by paper/tables/make_tables.py; do not edit.
% Inputs: live/live_metrics.csv, live/live_parity.csv
\begin{table*}[t]
\centering
\footnotesize
\setlength{\tabcolsep}{3pt}
\caption{Live HTTP replay (measured on loopback with an emulated network; RTT, bandwidth and object sizes are ASSUMED). Milliseconds of wait saved per view against the no-prefetch baseline: mean over citable passes (valid, and at most 1\% foreign I/O on the store device), with the envelope of the per-pass 95\% session-bootstrap CIs. Parity: sessions whose live counters differ from the simulator's.}
\label{tab:t7_live}
\begin{tabular}{llrrrrrrrr}
\hline
Store & Policy & Passes & Hit (\%) & Wait (ms) & Saved (ms) & Ideal saved & Wasted KiB & Bytes & Parity \\
\hline
\multicolumn{10}{l}{\emph{hdd object store}} \\
\hline
hdd & adaptive\_0.01 & 3 & $27.3$ & $109.1$ & $4.98$ {\scriptsize [$4.14$, $5.89$]} & $3.19$ & $94.6$ & $2.19\times$ & 0 \\
hdd & adaptive\_0.1@trans\_online & 3 & $29.2$ & $105.8$ & $8.27$ {\scriptsize [$7.09$, $9.73$]} & $6.12$ & $95.4$ & $2.20\times$ & 0 \\
hdd & feedback\_0.01\_eps0.005 & 3 & $27.9$ & $107.1$ & $6.99$ {\scriptsize [$5.95$, $8.19$]} & $4.34$ & $159.4$ & $3.00\times$ & 0 \\
hdd & fixed0 & 3 & $24.9$ & $114.1$ & $0.00$ {\scriptsize [$0.00$, $0.00$]} & $0.00$ & $0.0$ & $1.00\times$ & 0 \\
hdd & fixed4 & 3 & $28.4$ & $105.9$ & $8.18$ {\scriptsize [$6.91$, $9.71$]} & $5.09$ & $455.0$ & $6.71\times$ & 0 \\
hdd & fixed8 & 3 & $27.4$ & $106.9$ & $7.23$ {\scriptsize [$6.33$, $8.14$]} & $3.32$ & $1108.8$ & $14.91\times$ & 0 \\
\hline
\multicolumn{10}{l}{\emph{ssd object store}} \\
\hline
ssd & adaptive\_0.01 & 1 & $27.3$ & $101.8$ & $3.39$ {\scriptsize [$2.75$, $4.13$]} & $3.19$ & $94.6$ & $2.19\times$ & 0 \\
ssd & adaptive\_0.1@trans\_online & 2 & $29.2$ & $98.8$ & $6.36$ {\scriptsize [$5.34$, $7.51$]} & $6.12$ & $95.4$ & $2.20\times$ & 0 \\
ssd & feedback\_0.01\_eps0.005 & 2 & $27.9$ & $100.6$ & $4.58$ {\scriptsize [$3.72$, $5.67$]} & $4.34$ & $159.4$ & $3.00\times$ & 0 \\
ssd & fixed0 & 3 & $24.9$ & $105.2$ & $0.00$ {\scriptsize [$0.00$, $0.00$]} & $0.00$ & $0.0$ & $1.00\times$ & 0 \\
ssd & fixed4 & 3 & $28.4$ & $99.8$ & $5.40$ {\scriptsize [$4.45$, $6.53$]} & $5.09$ & $455.0$ & $6.71\times$ & 0 \\
ssd & fixed8 & 2 & $27.4$ & $101.8$ & $3.39$ {\scriptsize [$2.64$, $4.14$]} & $3.32$ & $1108.8$ & $14.91\times$ & 0 \\
\hline
\end{tabular}
\\[2pt]{\scriptsize 0 headline pass-policy row(s) failed the validity gate and 5 valid row(s) carried more than 1\% foreign I/O on the store device; both are excluded.}
\end{table*}


Every saving here is measured on loopback with an emulated \liveRttMs{} round
trip, a \liveMbps{} link and assumed object sizes; it is not a page-load
speedup of a real site. All \liveRows{} pass--policy rows passed the validity
gate, with \liveParityMismatches{} sessions whose live counters differ from the
simulator's (H1); \liveFlaggedRows{} SSD-store rows had over one percent foreign
I/O and are not cited.

\textbf{Prefetching lowered the wait.} With the hard-disk object store, all
\decTallyHTwoHddLowersTheMeanWait{} prefetching policies lowered the mean wait
against no prefetching in every repetition (H2; Table~\ref{tab:t7_live}). The
online-index controller (\polOnline) saved
\liveSavedHddAdaptiveZeroPointOneAtTransOnlineCI{} per view and \texttt{fixed4}
\liveSavedHddFixedFourCI, against a no-prefetch wait of \liveBaseWaitHdd; the
bracket is the envelope of the per-repetition session-bootstrap intervals. On
the SSD store \decTallyHTwoSsdLowersTheMeanWait{} of \decTallyHTwoSsd{}
policies lowered the wait; for \polAdaptive{} the effect is
\decLiveHTwoSsdAdaptiveZeroPointZeroOneVerdict, because only
\liveRepsSsdAdaptiveZeroPointZeroOne{} of its \liveHeadlineReps{} repetitions was
unflagged.

\textbf{The instant-prefetch model understated the hard-disk benefit.} The
model-ideal wait is zero on a hit and the round trip plus transfer on a miss.
On the hard-disk store the measured saving exceeded the model-ideal saving for
all \decTallyHThreeHddTheInstantPrefetchModelUnderstatesTheBenefit{} policies
(H3): \texttt{fixed8} saved \liveSavedHddFixedEight{} against a model-ideal
\liveIdealSavedHddFixedEight. The ideal model has no disk; a plausible reason is that
every demand miss on the hard disk also waits for a seek, which a prefetched
object avoids. On
the SSD store all \decTallyHThreeSsdConsistentWithTheModel{} policies were
consistent with the model, and the store changed the saving for
\decTallyHSevenTheStoreChangesTheDelta{} of \decTallyHSeven{} policies (H7).

\textbf{Bytes, scheduling and hints.} The savings cost bandwidth:
\texttt{fixed8} transferred \liveBytesXHddFixedEight{} the bytes of no
prefetching and saved less than \texttt{fixed4}; the two validated adaptive
controllers moved \liveBytesXHddAdaptiveZeroPointZeroOne{} and
\liveBytesXHddAdaptiveZeroPointOneAtTransOnline, and they are the only policies
no other policy beats on both saved wait and bytes (H4, descriptive;
Fig.~\ref{fig:f7}). Demand priority on the link matters for deep prefetching:
under a FIFO link, \texttt{fixed8}'s change in wait was
\decLiveHFiveSsdFixedEightFIFOCI, against \decLiveHFiveSsdFixedEightPriorityRepOneCI{}
with priority, so \decLiveHFiveSsdFixedEightVerdict{} (H5); for \texttt{fixed4}
there was \decLiveHFiveSsdFixedFourVerdict. Server-side read-ahead hints alone,
with no client prefetch, lowered the wait in one repetition on the hard disk
(H6; \decLiveHSixHddHintFixedEightCI{} for eight hinted objects).
